\icode{N24Z-MCH-SI-L3-261026-001}
\nceeID[3]{2024浙江1月}{8}
如图所示，小明取山泉水时发现水平细水管到水平地面的距离为水桶高的两倍，在地面上平移水桶，水恰好从桶口中心无阻挡地落到桶底边沿 \(A\)。已知桶高为 \(h\)，直径为 \(D\)，则水离开出水口的速度大小为\paren[]
	
\begin{tikzpicture}[font=\scriptsize]
	\draw[cyan] (0,0)node[above left,black]{\(A\)} parabola bend(2.4,3) (2.4,3);
	\draw[fill=brown!60] (2.4,2.95) rectangle(3,3.05);
	\draw[very thick] (0,0) rectangle(1.42,1.5);
	\draw[fill=gray!50] (.56,1.5)--(0.61,1.6)--(.81,1.6)--(.86,1.5);
	\foreach \x in {-.6,-.4,...,2}
		\draw[gray] (\x,-.2)--++(.2,.2);
	\draw[thick] (-.6,0)--(2,0);
	\node at(2.7,3.05)[above]{ 细水管};
\end{tikzpicture}
\begin{choices}
	\item \(\dfrac{D}{4} \sqrt{\dfrac{g}{h}}\)
	\item \(\dfrac{D}{2} \sqrt{\dfrac{g}{2h}}\)
	\item \(\dfrac{(\sqrt{2}+1)D}{2} \sqrt{\dfrac{g}{2h}}\)
	\item \((\sqrt{2}+1)D \sqrt{\dfrac{g}{2h}}\)
\end{choices}
\begin{answer}
	C\qquad 提示：水进入桶口到\(A\)的水平位移\(\frac{D}{2}\)，时间为\((\sqrt{2} - 1) \sqrt{\frac{2h}{g}}\)
\end{answer}